\section{What the Typology May and May Not Say}

\textbf{It may say} that canonical Scripture presents one Christ: the Word
made flesh, born at Bethlehem, laid where creatures feed, giving his flesh for
the world's life, and received sacramentally by the Church. Jerome, Ambrose,
and Chrysostom show that ancient Christians made substantial parts of this
connection. The image fittingly confesses the humility and generosity of the
Incarnation.

\textbf{It may not say} that the traditional gloss of Bethlehem proves a divine
code, that Luke secretly narrates the institution of the Eucharist, that
swaddling cloths are consecration linens, or that every Father read the manger
this way. It may not pass a pseudonymous homily off as Gregory Thaumaturgus or
an unattributed internet sentence off as Cyril of Alexandria. Resemblance
becomes responsible theology only when its evidentiary level is named.

\textbf{It must retain John 6's moral pressure.} The chapter is not a puzzle
whose reward is discovering Christmas symbols. Bread from heaven draws the
believer into Christ's self-gift, communion, abiding, resurrection, and the
unity of his Body. Chrysostom and Augustine alike refuse a reception that ends
at possession of a sacred object while leaving the communicant unchanged.

